% ==============================================================================
% compare_equipo_02.tex -- Algorithm Comparison (Problem Set 2: Predicting Poverty)
% Compilar:  cd 03_slides && latexmk -pdf compare_equipo_02.tex
% Las cifras vienen de numbers.tex (generado por 01_code/17_harold_diapositivas_figuras_y_cifras.R)
% y los puntajes de Kaggle de kaggle_scores.tex (se completan a mano).
% ==============================================================================
\def\deckname{Algorithm Comparison}
% ==============================================================================
% preamble.tex -- Estilo común de las dos presentaciones (Equipo 02)
% Compilar con pdflatex (sin paquetes de idioma ni fuentes especiales).
% ==============================================================================
\documentclass[aspectratio=169,9pt]{beamer}
\usepackage[utf8]{inputenc}
\usepackage[T1]{fontenc}
\usepackage{helvet}
\renewcommand{\familydefault}{\sfdefault}
\usepackage{booktabs, graphicx, xcolor, array, microtype, amsmath}
\graphicspath{{../02_outputs/figures/}}

% --- Colores (los mismos de las figuras) --------------------------------------
\definecolor{azul}{HTML}{1B3A6B}
\definecolor{naranja}{HTML}{E07B00}
\definecolor{verde}{HTML}{2A9D8F}
\definecolor{gris}{HTML}{5B6575}
\definecolor{fondo}{HTML}{F2F5FA}

% --- Tema ----------------------------------------------------------------------
\usetheme{default}
\setbeamertemplate{navigation symbols}{}
\setbeamercolor{normal text}{fg=black!85}
\setbeamercolor{structure}{fg=azul}
\setbeamercolor{block title}{bg=azul, fg=white}
\setbeamercolor{block body}{bg=fondo, fg=black!85}
\setbeamerfont{block title}{size=\small, series=\bfseries}
\setbeamertemplate{itemize item}{\color{naranja}\tiny$\blacksquare$}
\setbeamertemplate{itemize subitem}{\color{naranja}\tiny$\blacksquare$}
\setbeamertemplate{blocks}[rounded][shadow=false]
\setlength{\leftmargini}{1.1em}

\setbeamertemplate{frametitle}{%
  \vspace{1.5mm}
  \begin{minipage}[t]{\dimexpr\paperwidth-14mm\relax}
    \raggedright\color{azul}\bfseries\fontsize{12.5}{14}\selectfont\insertframetitle\par
    \vspace{0.8mm}{\color{naranja}\rule{12mm}{1.2pt}}
  \end{minipage}}
\setbeamertemplate{footline}{%
  \leavevmode\hbox{\hspace{7mm}%
    \begin{beamercolorbox}[wd=.80\paperwidth,ht=3.5ex,dp=1.2ex,left]{footline}
      {\color{gris}\scriptsize Equipo 02 · \deckname}
    \end{beamercolorbox}%
    \begin{beamercolorbox}[wd=.14\paperwidth,ht=3.5ex,dp=1.2ex,right]{footline}
      {\color{gris}\scriptsize\insertframenumber\,/\,\inserttotalframenumber}
    \end{beamercolorbox}}\vspace{1.5mm}}
\setbeamersize{text margin left=7mm, text margin right=7mm}

% --- Atajos --------------------------------------------------------------------
\newcommand{\pend}[1]{\textcolor{red}{\textbf{[#1]}}}
\newcommand{\estadistica}[2]{{\centering{\color{azul}\bfseries\fontsize{21}{23}\selectfont #1}\par{\color{gris}\footnotesize #2}\par}}
\newcommand{\figura}[2][.70\textheight]{\includegraphics[width=\linewidth,height=#1,keepaspectratio]{#2}}
\newcommand{\fuente}[1]{{\color{gris}\tiny #1}}
\newcommand{\resaltar}[1]{\textcolor{naranja}{\textbf{#1}}}

% --- Portada (se apaga con \portadafalse) ----------------------------------------------------
\newif\ifportada \portadatrue
\newcommand{\portada}[2]{%
  {\setbeamercolor{background canvas}{bg=azul}
  \begin{frame}[plain]
    \vspace*{9mm}
    {\color{white}\bfseries\fontsize{24}{27}\selectfont #1\par}
    \vspace{2mm}{\color{naranja}\rule{22mm}{2pt}}\par\vspace{4mm}
    {\color{white!90}\fontsize{13}{16}\selectfont #2\par}
    \vfill
    {\color{white}\bfseries\normalsize Equipo 02\par}
    \vspace{1mm}
    {\color{white!85}\small Jhon Sebasti\'an Vela Salcedo \ $\cdot$ \ Cesar Augusto Vallejo Quiroga \ $\cdot$ \ Harold Santiago Aguilera Mulfo\par}
    \vspace{4mm}
    {\color{white!70}\footnotesize Big Data y Machine Learning para Econom\'ia Aplicada $\cdot$ 2026-20 $\cdot$ Universidad de los Andes\par}
    \vspace{4mm}
  \end{frame}}}
% Generado por 01_code/17_harold_slide_assets.R; no editar a mano.
\newcommand{\nPobreTrain}{20,7}
\newcommand{\nNTrain}{164.960}
\newcommand{\nNTest}{66.168}
\newcommand{\nAUCRazon}{0,941}
\newcommand{\nCorrIng}{0,79}
\newcommand{\nAUCHacin}{0,696}
\newcommand{\nFoneA}{0,748}
\newcommand{\nFoneAD}{0,775}
\newcommand{\nFoneAll}{0,776}
\newcommand{\nFoneAB}{0,749}
\newcommand{\nFoneAC}{0,749}
\newcommand{\nFoneAblAtHalf}{0,767}
\newcommand{\nFoneSearchMin}{0,7729}
\newcommand{\nFoneSearchMax}{0,7770}
\newcommand{\nFoneRef}{0,7763}
\newcommand{\nNConfigs}{11}
\newcommand{\nShareSHAPD}{59}
\newcommand{\nShareSHAPTopThree}{43}
\newcommand{\nThrLGBM}{0,39}
\newcommand{\nFoneLGBM}{0,777}
\newcommand{\nPrecLGBM}{0,746}
\newcommand{\nRecLGBM}{0,811}
\newcommand{\nLeakLGBM}{25,4}
\newcommand{\nUnderLGBM}{18,9}
\newcommand{\nAUCLGBM}{0,958}
\newcommand{\nFoneLGBMsd}{0,0059}
\newcommand{\nPctPredPobre}{22,5}
\newcommand{\nPctRealPobre}{20,7}
\newcommand{\nFoneLGBMAtHalf}{0,768}
\newcommand{\nGainThr}{0,009}
\newcommand{\nUnderAtHalf}{25,9}
\newcommand{\nLeakAtHalf}{20,2}
\newcommand{\nThrUnderTen}{0,23}
\newcommand{\nLeakAtUnderTen}{34,9}
\newcommand{\nUnderCabecera}{18,8}
\newcommand{\nUnderRestorural}{19,9}
\newcommand{\nUnderJefaturafemenina}{17,0}
\newcommand{\nUnderJefaturamasculina}{20,6}
\newcommand{\nUnderPropia}{21,3}
\newcommand{\nUnderArriendo}{20,0}
\newcommand{\nUnderUsufructouotra}{14,9}
\newcommand{\nFoneSurveyMin}{0,585}
\newcommand{\nUnderSurveyMin}{34}
\newcommand{\nLeakSurveyMin}{48}
\newcommand{\nNVarsSurveyMin}{54}
\newcommand{\nFoneSurveyShort}{0,675}
\newcommand{\nUnderSurveyShort}{26}
\newcommand{\nLeakSurveyShort}{38}
\newcommand{\nNVarsSurveyShort}{75}
\newcommand{\nFoneSurveyFull}{0,777}
\newcommand{\nUnderSurveyFull}{19}
\newcommand{\nLeakSurveyFull}{25}
\newcommand{\nNVarsSurveyFull}{149}
\newcommand{\nSPBoostingLightGBM}{0,777}
\newcommand{\nSPRandomForest}{0,774}
\newcommand{\nSPLogit}{0,771}
\newcommand{\nSPElasticNet}{0,771}
\newcommand{\nSPLPM}{0,770}
\newcommand{\nSPCART}{0,769}
\newcommand{\nSPNaiveBayes}{0,758}
\newcommand{\nSPRange}{0,008}
\newcommand{\nGainIncome}{0,027}
\newcommand{\nGainAlgo}{0,006}
\newcommand{\nGainHyper}{0,0041}
\newcommand{\nRondasLGBM}{954}
\newcommand{\nNVarsLGBM}{149}
\newcommand{\nHojasLGBM}{15}
\newcommand{\nPublicSD}{0,0057}
\newcommand{\nDropLogitRegions}{0,0010}
\newcommand{\nDropLGBMRegions}{0,0031}
\newcommand{\nGapLogit}{0,0006}
\newcommand{\nDisagreeShare}{2,3}
\newcommand{\nDisagreePoor}{37}
\newcommand{\nFoneLGBMAll}{0,777}
\newcommand{\nPrecLGBMAll}{0,746}
\newcommand{\nRecLGBMAll}{0,811}
\newcommand{\nFoneXGB}{0,775}
\newcommand{\nPrecXGB}{0,749}
\newcommand{\nRecXGB}{0,803}
\newcommand{\nFoneRF}{0,774}
\newcommand{\nPrecRF}{0,748}
\newcommand{\nRecRF}{0,801}
\newcommand{\nFoneEN}{0,772}
\newcommand{\nPrecEN}{0,730}
\newcommand{\nRecEN}{0,819}
\newcommand{\nFoneLogitRob}{0,771}
\newcommand{\nPrecLogitRob}{0,739}
\newcommand{\nRecLogitRob}{0,806}
\newcommand{\nFoneLogitFortyVars}{0,771}
\newcommand{\nPrecLogitFortyVars}{0,738}
\newcommand{\nRecLogitFortyVars}{0,807}
\newcommand{\nFoneProbitRob}{0,771}
\newcommand{\nPrecProbitRob}{0,739}
\newcommand{\nRecProbitRob}{0,807}
\newcommand{\nFoneLPM}{0,770}
\newcommand{\nPrecLPM}{0,728}
\newcommand{\nRecLPM}{0,816}
\newcommand{\nFoneCART}{0,769}
\newcommand{\nPrecCART}{0,755}
\newcommand{\nRecCART}{0,783}
\newcommand{\nFoneNB}{0,758}
\newcommand{\nPrecNB}{0,694}
\newcommand{\nRecNB}{0,834}
\newcommand{\nMtryRF}{14}
\newcommand{\nNodeRF}{150}

% Puntajes de Kaggle (F1 de la clase pobre) que citan las diapositivas: público (20 % del test) y privado (80 %).
\newcommand{\kgLGBM}{0,78}   \newcommand{\kgpLGBM}{0,77}   % LightGBM, 149 variables (harold_lightgbm_fe_leaves_15_rounds_954.csv)
\newcommand{\kgRF}{0,77}     \newcommand{\kgpRF}{0,77}     % Random Forest (harold_random_forest_fe_mtry_14_node_150.csv)

% Generado por 01_code/18_harold_diapositivas_envios_equipo.R; no editar a mano.
\newcommand{\KgEnvios}{22}
\newcommand{\KgPrivMasSetentaYSiete}{10}
\newcommand{\KgPubMasSetentaYSiete}{10}
\newcommand{\KgPrivMax}{0,78}
\newcommand{\KgPubMax}{0,78}
\newcommand{\KgNPrivMax}{3}
\newcommand{\KgEnviosHarold}{6}
\newcommand{\KgEnviosJhon}{9}
\newcommand{\KgEnviosCesar}{7}
\newcommand{\NModelosCerca}{siete}
\newcommand{\NModelosTotal}{ocho}
\newcommand{\SPRangeCinco}{0,007}
\newcommand{\JhFoneLGBMocho}{0,777}
\newcommand{\JhFoneRFveinte}{0,778}
\newcommand{\JhPrecRFveinte}{0,740}
\newcommand{\JhRecRFveinte}{0,820}
\newcommand{\JhFoneLGBMsin}{0,765}
\newcommand{\JhFoneXGBsin}{0,761}
\newcommand{\JhFoneRFsin}{0,721}
\newcommand{\JhFoneXGBocho}{0,777}
\newcommand{\JhFoneLGBMmil}{0,773}
\newcommand{\JhFoneCorta}{0,706}
\newcommand{\JhFoneLogitCorta}{0,689}
   % cifras de los envíos de todo el equipo (script 18)

\begin{document}

\ifportada
\portada{Problem Set 2: Predicting Poverty}{Algorithm Comparison: qu\'e diferencia a los algoritmos y cu\'al recomendar\'iamos}
\fi

% ------------------------------------------------------------------------------
% 1. Resumen del mejor modelo (diapositiva obligatoria)
% ------------------------------------------------------------------------------
\begin{frame}{Best Model Overview: LightGBM gana por poco; los datos explican más que el algoritmo}
  \begin{columns}[T]
    \column{.50\textwidth}
    \begin{block}{Especificación del mejor modelo}
      \footnotesize
      \begin{itemize}\setlength{\itemsep}{1pt}
        \item \textbf{Algoritmo:} LightGBM, clasificación binaria (boosting de árboles)
        \item \textbf{Variables:} \nNVarsLGBM{} a nivel de hogar, con el \resaltar{ingreso imputado} por persona
        \item \textbf{Hiperparámetros:} \nHojasLGBM{} hojas, \nRondasLGBM{} rondas, tasa 0,03, L2 de 50
        \item \textbf{Umbral:} \nThrLGBM{} (maximiza el F1); validación en 5 pliegues, sin Bogotá
      \end{itemize}
    \end{block}
    \column{.22\textwidth}
    \estadistica{\kgLGBM}{F1 público\\en Kaggle}
    \vspace{2mm}
    \estadistica{\nFoneLGBM}{F1 en validación\\cruzada}
    \column{.22\textwidth}
    \estadistica{\KgPrivMax}{F1 privado, mejor\\del equipo}
    \vspace{2mm}
    \estadistica{\SPRangeCinco}{rango de F1 de cinco\\algoritmos, mismas variables}
  \end{columns}
  \begin{block}{Mensajes principales}
    \footnotesize
    \begin{enumerate}\setlength{\itemsep}{1pt}
      \item Con las \resaltar{mismas 40 variables}, cinco de los seis algoritmos quedan a \SPRangeCinco{} de F1; solo Naive Bayes se aleja (\nSPNaiveBayes{}).
      \item Lo que más mueve el F1 es el ingreso imputado (+\nGainIncome{}), no el algoritmo (+\nGainAlgo{} de Logit a LightGBM).
      \item Random Forest y Elastic Net, los contendientes más cercanos, pierden menos de 0,006. El mejor privado (\KgPrivMax) lo logran modelos con \resaltar{8 a 20 variables} y la misma idea de ingreso imputado: todo cabe en el ruido de la muestra pública.
    \end{enumerate}
  \end{block}
\end{frame}

% ------------------------------------------------------------------------------
% Contenido (después de la diapositiva obligatoria de resumen)
% ------------------------------------------------------------------------------
\ifportada
\begin{frame}{Contenido: de lo que dicen los datos a qué modelo recomendar}
  \renewcommand{\arraystretch}{1.9}
  \centering\normalsize
  \begin{tabular}{@{}r@{\hspace{4mm}}p{3.4cm}p{9.6cm}@{}}
    {\color{naranja}\bfseries 2} & \textbf{Resumen} & Cuál es el mejor modelo y qué dice la comparación en tres mensajes\\
    {\color{naranja}\bfseries \getpagerefnumber{p:datos}--\the\numexpr\getpagerefnumber{p:comp}-1\relax} & \textbf{Datos} & Por qué un modelo lineal puede competir con los árboles y cómo se compara de forma justa\\
    {\color{naranja}\bfseries \getpagerefnumber{p:comp}--\the\numexpr\getpagerefnumber{p:concl}-1\relax} & \textbf{Comparación} & Los algoritmos exigidos, los dos contendientes más cercanos y qué explica la brecha\\
    {\color{naranja}\bfseries \getpagerefnumber{p:concl}} & \textbf{Conclusión} & Complejidad, interpretabilidad y desempeño: qué modelo recomendaríamos a una agencia
  \end{tabular}
\end{frame}
\fi

% ------------------------------------------------------------------------------
% DATOS
% ------------------------------------------------------------------------------
\begin{frame}{Datos: la variable clave es casi logística, por eso un modelo lineal compite con los árboles}
  \label{p:datos}
  \begin{columns}[T]
    \column{.64\textwidth}
    \figura{harold_slide_nonlinearity.png}
    \column{.34\textwidth}
    \small
    \begin{itemize}\setlength{\itemsep}{3pt}
      \item La probabilidad de pobreza depende sobre todo de la \resaltar{razón ingreso imputado / línea}, con forma de S.
      \item Un logit con splines sobre esa variable queda casi pegado al LightGBM.
      \item Los árboles ganan donde hay interacciones; después de construir esa razón quedan pocas.
      \item El desbalance (\nPobreTrain{}\,\% pobres) se maneja con el \resaltar{umbral}: de 0,5 a \nThrLGBM{} el F1 sube \nGainThr{}.
    \end{itemize}
  \end{columns}
\end{frame}

\begin{frame}{Datos: una comparación justa fija predictores, pliegues y regla de umbral}
  \begin{columns}[T]
    \column{.50\textwidth}
    \small
    \begin{itemize}\setlength{\itemsep}{4pt}
      \item \textbf{Mismas 40 variables} para los seis algoritmos, las de mayor importancia SHAP (selección con un modelo auxiliar).
      \item \textbf{Mismos 5 pliegues} de hogares y la misma métrica: F1 de la clase pobre, sin Bogotá (no existe en test).
      \item \textbf{Mismo umbral:} el que maximiza el F1 con las predicciones fuera de muestra; cada algoritmo ajusta su umbral, no el puntaje de Kaggle.
      \item \textbf{Mismo preprocesamiento:} los modelos lineales usan log1p, winsorización y dummies; los árboles, las variables sin transformar.
    \end{itemize}
    \column{.46\textwidth}
    \begin{block}{Qué cambia entre algoritmos}
      \footnotesize
      Solo la forma funcional y la regularización: lineal sin penalizar (LPM, Logit), penalizado (Elastic Net), promedio de árboles (Random Forest), árboles en secuencia (LightGBM) y probabilidades condicionales independientes (Naive Bayes).
    \end{block}
    \vspace{1mm}
    \fuente{Las versiones enviadas a Kaggle usan más variables (LightGBM, 149; Elastic Net, 102 columnas con splines; Logit, 59). Entre versiones de distinto protocolo solo se compara el puntaje de Kaggle.}
  \end{columns}
\end{frame}

% ------------------------------------------------------------------------------
% COMPARACIÓN
% ------------------------------------------------------------------------------
\begin{frame}{Resumen: la mejor versión de cada uno de los seis algoritmos}
  \label{p:comp}
  \centering
  \small
  \resizebox{.97\linewidth}{!}{\begin{tabular}{llcccccc}
\toprule
 & & \multicolumn{3}{c}{\textbf{Validación cruzada}} & \multicolumn{3}{c}{\textbf{Kaggle}} \\
\cmidrule(lr){3-5}\cmidrule(lr){6-8}
\textbf{Algoritmo} & \textbf{Versión de referencia} & \textbf{F1} & \textbf{Prec.} & \textbf{Cob.} & \textbf{Envíos} & \textbf{Máx. público} & \textbf{Máx. privado} \\
\midrule
Boosting & LightGBM, 149 var. & 0,777 & 0,746 & 0,811 & 7 & 0,78 & 0,78 \\
Random Forest & ranger, 40 var. & 0,774 & 0,748 & 0,801 & 4 & 0,77 & 0,78 \\
Elastic Net & glmnet, splines & 0,772 & 0,730 & 0,819 & 5 & 0,77 & 0,77 \\
Logit & regularizado, 59 var. & 0,771 & 0,739 & 0,806 & 4 & 0,77 & 0,77 \\
LPM & MCO, 40 var. & 0,770 & 0,728 & 0,816 & 1 & 0,74 & 0,75 \\
Naive Bayes & deciles, 40 var. & 0,758 & 0,694 & 0,834 & 1 & 0,64 & 0,65 \\
\bottomrule
\end{tabular}
}\\[3mm]
  \footnotesize\raggedright
  Validación cruzada de la versión de referencia: 5 pliegues, sin Bogotá, umbral que maximiza el F1. Kaggle: mejor puntaje público y mejor privado entre los envíos de cada algoritmo.
\end{frame}

\begin{frame}{Con cada algoritmo en su mejor versión, \NModelosCerca{} de \NModelosTotal{} modelos quedan a menos de un punto}
  \begin{columns}[T]
    \column{.64\textwidth}
    \figura{harold_slide_model_f1_compare.png}
    \column{.34\textwidth}
    \small
    \begin{itemize}\setlength{\itemsep}{3pt}
      \item Las barras (±1 DE entre pliegues) se traslapan: la mayoría de las diferencias es del tamaño del ruido.
      \item Cada modelo usa sus propias variables (40, 59 o \nNVarsLGBM), así que esta comparación aún no es justa.
      \item Solo Naive Bayes queda claramente atrás: supone variables independientes, y aquí no lo son.
    \end{itemize}
  \end{columns}
\end{frame}

\begin{frame}{Con las mismas variables, cinco de seis algoritmos quedan a menos de un punto de F1}
  \begin{columns}[T]
    \column{.64\textwidth}
    \figura{harold_slide_same_predictors_compare.png}
    \column{.34\textwidth}
    \small
    \begin{itemize}\setlength{\itemsep}{3pt}
      \item El Elastic Net con 40 variables queda en \nSPElasticNet{}: la penalización casi no actúa porque hay 165 mil hogares para pocos parámetros.
      \item \textbf{Naive Bayes} rinde menos: supone independencia entre variables que están muy correlacionadas.
    \end{itemize}
  \end{columns}
\end{frame}

\begin{frame}{Contendiente 1, Random Forest: promedia árboles y, con 20 variables, llega a \KgPrivMax{} privado}
  \begin{columns}[T]
    \column{.41\textwidth}
    \small
    \begin{itemize}\setlength{\itemsep}{3pt}
      \item \textbf{Bosque:} cientos de árboles profundos y decorrelacionados promedian el ruido.
      \item \textbf{Boosting:} cada árbol corrige el residuo del anterior y reduce el sesgo.
      \item Ambos priorizan la razón ingreso imputado / línea y variantes de ella.
    \end{itemize}
    \vspace{1mm}
    \begin{block}{Dos versiones del mismo algoritmo}
      \footnotesize\centering
      \resizebox{\linewidth}{!}{\begin{tabular}{lcc}
        \toprule
         & \textbf{Versión A} & \textbf{Versión B} \\
        \midrule
        Variables & 40 & 20 \\
        Árboles; hoja mínima & 500; \nNodeRF & 360; 156 \\
        F1 (validación cruzada) & \nSPRandomForest & \JhFoneRFveinte \\
        Precisión / cobertura & \nPrecRF{} / \nRecRF & \JhPrecRFveinte{} / \JhRecRFveinte \\
        Kaggle público / privado & \kgRF{} / \kgpRF & 0,77 / 0,78 \\
        \bottomrule
      \end{tabular}}\\[1mm]
      \raggedright Referencia, LightGBM de 149 variables: F1 \nSPBoostingLightGBM{}, precisión \nPrecLGBM{}, cobertura \nRecLGBM.
    \end{block}
    \column{.57\textwidth}
    \centering\small\textbf{Posición de cada variable en la importancia}\\[2mm]
    \footnotesize
    \resizebox{\linewidth}{!}{\begin{tabular}{lccc}
\toprule
\textbf{Variable} & \textbf{LightGBM} & \textbf{Random Forest} & \textbf{Elastic Net} \\
\midrule
Log razón ingreso imputado / línea & 1 & 2 & 4 \\
Razón ingreso imputado / línea & 2 & 1 & 2 \\
Arriendo imputado cubre la línea & 39 & -- & 1 \\
Brecha de ingreso imputado (% de la línea) & 3 & 3 & -- \\
Ingreso laboral requerido & 11 & 9 & 3 \\
Personas con subsidios & 4 & 4 & 12 \\
Arriendo imputado / línea & 5 & 5 & 9 \\
Ingreso imputado per cápita & 33 & 29 & 5 \\
Ingreso imputado del hogar & 6 & 6 & 6 \\
\bottomrule
\end{tabular}
}\\[2mm]
    \fuente{LightGBM y Random Forest: valor SHAP absoluto medio; Elastic Net: mayor coeficiente estandarizado en valor absoluto. ``--'': fuera del modelo. El F1 de las dos versiones viene de protocolos de validación distintos.}
  \end{columns}
\end{frame}

\begin{frame}{Contendiente 2, Elastic Net: casi los mismos hogares, con coeficientes que se pueden leer}
  \begin{columns}[T]
    \column{.40\textwidth}
    \figura{harold_slide_en_coefs.png}
    \column{.56\textwidth}
    \centering\small\textbf{LightGBM frente a Elastic Net: hogares clasificados distinto}\\[2mm]
    \footnotesize
    \resizebox{\linewidth}{!}{\begin{tabular}{lrrr}
\toprule
\textbf{Clasificación} & \textbf{Hogares} & \textbf{\% del total} & \textbf{\% realmente pobres} \\
\midrule
Ambos: pobre & 33.514 & 21,7 & 76 \\
Ambos: no pobre & 117.275 & 76,0 & 4 \\
Solo LightGBM: pobre & 1.227 & 0,8 & 42 \\
Solo Elastic Net: pobre & 2.377 & 1,5 & 34 \\
\bottomrule
\end{tabular}
}\\[2mm]
    \raggedright\small
    \begin{itemize}\setlength{\itemsep}{2pt}
      \item Los modelos discrepan en \nDisagreeShare{}\,\% de los hogares y en ellos solo \nDisagreePoor{}\,\% son pobres: son casos de frontera.
      \item Los coeficientes del Elastic Net recuperan lo mismo que el SHAP: dominan la razón ingreso imputado y el arriendo imputado.
    \end{itemize}
  \end{columns}
\end{frame}

\begin{frame}{Qué explica la brecha: ingreso imputado, umbral y luego el algoritmo}
  \begin{columns}[T]
    \column{.66\textwidth}
    \figura{harold_slide_what_moves_f1.png}
    \column{.32\textwidth}
    \small
    \begin{itemize}\setlength{\itemsep}{3pt}
      \item Pasar de las variables base a las de ingreso imputado vale \nGainIncome{}: cuatro veces más que cambiar de algoritmo.
      \item Elegir el \resaltar{umbral} sobre las probabilidades vale más que el algoritmo.
      \item Afinar los hiperparámetros mueve el F1 \nGainHyper{} entre la mejor y la peor de 11 configuraciones.
      \item {\footnotesize\textbf{Lo replica otro modelo del equipo:} sin las variables de la primera etapa su F1 baja de $\approx$0,777 a \JhFoneLGBMsin{} (LightGBM), \JhFoneXGBsin{} (XGBoost) y \JhFoneRFsin{} (RF); con 1.152 variables de persona y sin primera etapa llega a \JhFoneLGBMmil.}
    \end{itemize}
  \end{columns}
\end{frame}

\begin{frame}{Más allá del Kaggle: el logit regularizado se parece al LightGBM y es más estable}
  \begin{columns}[T]
    \column{.56\textwidth}
    \centering\small\textbf{F1 en hogares nuevos y en regiones nuevas}\\[2mm]
    \footnotesize
    \resizebox{\linewidth}{!}{\begin{tabular}{lcccc}
\toprule
\textbf{Modelo} & \textbf{F1 dentro} & \textbf{F1 fuera} & \textbf{F1 regiones nuevas} & \textbf{Caída} \\
\midrule
Logit regularizado & 0,7719 & 0,7713 & 0,7703 & 0,0010 \\
LightGBM & -- & 0,7772 & 0,7741 & 0,0031 \\
\bottomrule
\end{tabular}
}\\[2mm]
    \fuente{Regiones nuevas: en cada pliegue se dejan fuera dominios geográficos completos.}
    \column{.40\textwidth}
    \small
    \begin{itemize}\setlength{\itemsep}{3pt}
      \item Entre logit regularizado (\nFoneLogitRob) y LightGBM (\nFoneLGBM) hay 0,006 de F1.
      \item En una muestra del tamaño del leaderboard público, la desviación estándar del F1 es \nPublicSD{}: la diferencia es del tamaño del ruido.
      \item Dentro y fuera de muestra el logit casi no cambia (brecha \nGapLogit); con regiones nuevas, cae \nDropLogitRegions{} contra \nDropLGBMRegions{} del boosting.
    \end{itemize}
  \end{columns}
\end{frame}

% ------------------------------------------------------------------------------
% CONCLUSIÓN
% ------------------------------------------------------------------------------
\begin{frame}{Conclusión: el mejor puntaje público no tiene por qué ser el modelo que recomendaríamos}
  \label{p:concl}
  \small
  \centering
  \resizebox{.95\linewidth}{!}{%
  \begin{tabular}{p{3.3cm}p{1.8cm}p{2.0cm}p{3.4cm}p{3.2cm}}
    \toprule
    \textbf{Modelo} & \textbf{F1 (CV)} & \textbf{Variables} & \textbf{Interpretación} & \textbf{Caída, regiones nuevas} \\
    \midrule
    LightGBM & \nFoneLGBM & \nNVarsLGBM & SHAP, no coeficientes & cae \nDropLGBMRegions \\
    LightGBM, 8 variables & \JhFoneLGBMocho\textsuperscript{*} & 8 & SHAP; pocas variables & no medida \\
    Random Forest & \nFoneRF & 40 & importancia global & no medida \\
    Logit regularizado & \nFoneLogitRob & 59 & coeficientes y signos & cae \nDropLogitRegions \\
    \bottomrule
  \end{tabular}}\\[0.5mm]
  {\scriptsize\color{gris} * Validación propia de este modelo (anidada, sin Bogotá); no es comparable con las demás filas.}
  \vspace{2mm}
  \begin{columns}[T]
    \column{.48\textwidth}
    \begin{block}{Qué dice la comparación}
      \small
      La complejidad compra \resaltar{menos de un punto de F1}: con 8 variables, un modelo logra el mejor privado (\KgPrivMax). La información (el ingreso imputado) explica casi todo y es lo difícil de desplegar: necesita el módulo laboral de cada persona.
    \end{block}
    \column{.48\textwidth}
    \begin{block}{Recomendación}
      \small
      Para una agencia: el \resaltar{logit regularizado}, auditable y estable entre regiones, con el LightGBM como control. Si se busca el máximo puntaje, LightGBM; la elección depende de cuánto valore la agencia explicar cada exclusión.
    \end{block}
  \end{columns}
\end{frame}

\end{document}
